\subsection{CALCULATION OF A DETERMINANT}

Lemma 1. Let A be a commutative ring and M a free A-module with a basis \((e_j)_{j \in J}\) , where the indexing set \(J\) is totally ordered. For every integer \(p \leqslant \operatorname{Card}(J)\) , every alternating \(p\) -linear function \(f: \mathbf{M}^p \to \mathbf{N}\) (where \(N\) is an A-module) and every family of \(p\) elements \(x_i = \sum_{j \in J} \xi_{ji} e_j\) of \(M\) ( \(1 \leqslant i \leqslant p\) ),

\[
\begin{array}{r l} f (x _ {1}, x _ {2}, \dots , x _ {p}) & = \sum_ {j _ {1} <   j _ {2} <   \dots <   j _ {p}} \left(\sum_ {\sigma \in \mathfrak {S} _ {p}} \varepsilon_ {\sigma}. \xi_ {j _ {\sigma (1)}, 1} \xi_ {j _ {\sigma (2)}, 2} \dots \xi_ {j _ {\sigma (p)}, p}\right) f (e _ {j _ {1}}, \dots , e _ {j _ {p}}) \end{array}\tag{10}
\]

where \((j_k)_{1\leqslant k\leqslant p}\) runs through the set of strictly increasing sequences of \(p\) elements of J.

Now

\[
f (x _ {1}, \dots , x _ {p}) = \sum_ {(j _ {k})} \xi_ {j _ {1}, 1} \xi_ {j _ {2}, 2} \dots \xi_ {j _ {p}, p} f (e _ {j _ {1}}, e _ {j _ {2}}, \dots , e _ {j _ {p}})
\]

where \((j_k)_{1 \leqslant k \leqslant p}\) runs through all the sequences of \(p\) elements of \(J\) ; it then suffices to apply Corollary 1 to Proposition 7 of §7, no. 4 to \(f\) .

In particular, if J is finite and has n elements and \(x_{i} = \sum_{j \in J} \xi_{ji} e_{j} (1 \leqslant i \leqslant n)\) are n elements of M, then

\[
\begin{array}{l} \text {(11)} \quad x _ {1} \wedge x _ {2} \wedge \dots \wedge x _ {n} \\ = \left(\sum_ {\sigma \in \mathfrak {S} _ {n}} \varepsilon_ {\sigma} \xi_ {j _ {\sigma (1)}, 1} \xi_ {j _ {\sigma (2)}, 2} \dots \xi_ {j _ {\sigma (n)}, n}\right) e _ {j _ {1}} \wedge e _ {j _ {2}} \wedge \dots \wedge e _ {j _ {n}} \end{array}
\]

where \((j_k)_{1\leqslant k\leqslant n}\) is the unique sequence of the \(n\) elements of J arranged in increasing order, whence

\[
\det (x _ {1}, x _ {2}, \dots , x _ {n}) = \sum_ {\sigma \in \mathfrak {S} _ {n}} \varepsilon_ {\sigma}. \xi_ {j _ {\sigma (1)}, 1} \xi_ {j _ {\sigma (2)}, 2} \dots \xi_ {j _ {\sigma (n)}, n}.\tag{12}
\]

With the notation of Lemma 1, comparing formulae (10) and (12) gives

\[
x _ {1} \wedge x _ {2} \wedge \dots \wedge x _ {p} = \sum_ {\mathrm{H} \in \mathfrak {F} _ {p} (J)} \det (x _ {\mathrm{H}, 1}, x _ {\mathrm{H}, 2}, \dots , x _ {\mathrm{H}, p}) e _ {\mathrm{H}}\tag{13}
\]

where \(\mathfrak{F}_{p}(\mathrm{J})\) is the set of subsets of J with p elements and, for every subset \(\mathrm{H}\in\mathfrak{F}_{p}(\mathrm{J})\) , we write \(x_{H,i}=\sum_{j\in H}\xi_{ji}e_{j}\) and \(e_{H}=e_{j_{1}}\wedge e_{j_{2}}\wedge\cdots\wedge e_{j_{p}},(j_{k})_{1\leqslant k\leqslant p}\) being the sequence of elements of H arranged in increasing order, it being understood that \(\det(x_{\mathrm{H},1},\ldots,x_{\mathrm{H},p})\) is taken with respect to the basis \((e_{jk})_{1\leqslant k\leqslant p}\) .

PROPOSITION 7. Let I be a finite set and \(X = (\xi_{ji})_{(j, i) \in I \times I}\) a square matrix of type (I, I) over a commutative ring A. Then

\[
\det X = \sum_ {\sigma \in \mathfrak {S} _ {I}} \varepsilon_ {\sigma} \left(\prod_ {i \in I} \xi_ {\sigma (i), i}\right)\tag{14}
\]

where \(\sigma\) runs through the group \(\mathfrak{S}_{\mathbf{I}}\) of permutations of \(\mathbf{I}\) and \(\varepsilon_{\sigma}\) is the signature of \(\sigma\) (I, § 5, no. 7).

Attention may be confined to the case where \(I = [1, n] \subset N\) and it then suffices to apply formula (12), where \((e_{i})_{1 \leqslant i \leqslant n}\) is the canonical basis of \(A^{n}\) and the \(x_{i}\) are the columns of X (cf.~no. 3, formula (6)).

In particular, for the determinant of a matrix of order 3

\[
X = \left( \begin{array}{c c c} \xi_ {1 1} & \xi_ {1 2} & \xi_ {1 3} \\ \xi_ {2 1} & \xi_ {2 2} & \xi_ {2 3} \\ \xi_ {3 1} & \xi_ {3 2} & \xi_ {3 3} \end{array} \right)
\]

we have

\[
\begin{array}{r l} \det (X) = & \xi_ {1 1} \xi_ {2 2} \xi_ {3 3} + \xi_ {1 2} \xi_ {2 3} \xi_ {3 1} + \xi_ {2 1} \xi_ {3 2} \xi_ {1 3} \\ & - \xi_ {1 3} \xi_ {2 2} \xi_ {3 1} - \xi_ {1 2} \xi_ {2 1} \xi_ {3 3} - \xi_ {1 1} \xi_ {2 3} \xi_ {3 2}. \end{array}
\]

PROPOSITION 8. For every square matrix X over a commutative ring, the determinant of the transpose matrix \(^{t}X\) is equal to the determinant of X.

Suppose that X is of type (I, I). For every ordered pair of permutations \(\sigma\) , \(\tau\) of \(S_{I}\) , (since multiplication is commutative)

\[
\prod_ {i \in I} \xi_ {\sigma (i), i} = \prod_ {j \in I} \xi_ {\sigma (\tau (j)), \tau (j)}.
\]

In particular take \(\tau = \sigma^{-1}\) ; using the fact that \(\varepsilon_{\sigma^{-1}} = \varepsilon_{\sigma}\) , it is seen that

\[
\det X = \sum_ {\sigma \in \mathfrak {S} _ {I}} \varepsilon_ {\sigma}. \left(\prod_ {i \in I} \xi_ {i, \sigma (i)}\right),\tag{15}
\]

which proves the proposition.

COROLLARY 1. For \(n\) vectors \(x_1, \ldots, x_n\) of \(A^n\) , let \(Y(x_1, \ldots, x_n)\) denote the square matrix of order \(n\) whose \(i\) -th row is \(x_i\) , for \(1 \leqslant i \leqslant n\) . Then the mapping

\[
(x _ {1}, \dots , x _ {n}) \mapsto \det (Y (x _ {1}, \dots , x _ {n}))
\]

from \((\mathbf{A}^n)^n\) to \(\mathbf{A}\) is alternating and \(n\) -linear.

COROLLARY 2. For a square matrix \(X\) of finite order over a commutative ring \(A\) the following conditions are equivalent:

\begin{enumerate}
\def\labelenumi{(\roman{enumi})}
\item
  the rows of \(X\) are linearly independent;
\item
  the columns of \(X\) are linearly independent;
\item
  det X is not a divisor of zero in A.
\end{enumerate}

This follows immediately from no. 3, Proposition 6 and Proposition 8 above.

COROLLARY 3. Let \(u\) be an endomorphism of a finite-dimensional free A-module \(M\) and \({}^t u\) the transpose endomorphism of the dual \(M^*\) (II, §2, no. 5, Definition 5); then

\[
\det (t u) = \det (u).\tag{16}
\]

If X is the matrix of u with respect to a basis of M, \({}^{t}X\) is the matrix of \({}^{t}u\) with respect to the dual basis (II, § 10, no. 4, Proposition 3); as \(\det(u) = \det(X)\) and \(\det(^{t}u) = \det(^{t}X)\) , the conclusion follows from Proposition 8.

